\vspace{-10pt}
\section{Introduction}



Due to the ever-increasing demand for computation from neural network workloads, various deep learning systems (DLS), e.g., NVIDIA GPU with Tensor Core~\cite{NVTensorCore}, Google TPU~\cite{Jouppi17ISCA}, GraphCore IPU~\cite{graphcoreipu}, and Cambricon MLU~\cite{MLU}, have been deployed in data centers of cloud and internet service companies such as Microsoft~\cite{Microsoftcloud}, Google~\cite{Googlecloud}, and Amazon~\cite{Amazoncloud}. To fully exploit various DLS, it is required to develop multiple high-performance tensor programs, i.e., low-level implementations of tensor operators in deep learning, for different platforms, which are notoriously challenging because of the complicated architecture and programming models. Ideally, if we only need to write one copy of a program, and can run it on different platforms, that is, “Write Once, Run Anywhere”, it is feasible to address the programming problem in data centers with heterogeneous DLS.


\begin{figure}[t]
    \vspace{-0pt}
    \centering
    \includegraphics[width=1.0\linewidth]{figures/fig1.pdf}
    %\vspace{-20pt}
    \caption{\small Comparing \name{} to existing transcompilation techniques in terms of (a) scalability-accuracy axis and (b) human efforts}
    \label{fig:intro_comparison}
    %\vspace{-10pt}
\end{figure}

One of the most attractive solutions to achieve “Write Once, Run Anywhere” on heterogeneous platforms is source-to-source compilation, i.e., transcompilation, which automatically compiles the code written with one high-level language to another. Existing transcompilation efforts fall into three categories, i.e., rule-based, symbolic synthesis, and data-driven approaches. The rule-based approaches require experts to manually define a set of transformation rules, typically applied to Abstract Syntax Tree (AST), between different programming languages, and then employ pattern matching to parse the input source programs~\cite{johnson2022martini,10.1145/3554736}. For example, a CUDA-to-FPGA translator, i.e., FCUDA, defines a set of transformation rules for data communication, compute optimization, and parallelism mapping~\cite{PAPAKONSTANTINOU13TECS}. The symbolic synthesis approaches generate semantic-preserving target code from either domain-specific languages~\cite{Phitchaya14PLDI} or input/output examples~\cite{Jackson22PLDI}. Since they typically rely on a costly search-based SMT (Satisfiability Modulo Theories) solver~\cite{Z3Solver}, it is hard to scale to large-size general-purpose programs. The data-driven approaches, which typically train neural networks from a large amount of source code for code generation, have emerged recently~\cite{bahdanau2014neural,wu2016google,Shah18NIPS}. Notable examples include TransCoder~\cite{Roziere20NIPS},  StarCoder~\cite{li2023starcoder}, GPT-4~\cite{achiam2023gpt}, and OpenAI-o1~\cite{openaio1} which automatically translate programs written in different high-level languages such as C++, Java, and Python.




However, the above transcompilation techniques are not applicable to DLS due to their complicated architecture and programming models. Regarding the rule-based approaches, the huge architectural discrepancy between different platforms makes it infeasible to manually define efficient transformation rules. Regarding the symbolic synthesis approaches, in addition to their limited scalability, it is infeasible to handle different parallel semantics, e.g., SIMT in CUDA C and SIMD in BANG C~\cite{bangc}, for the SMT solver. Meanwhile, it also requires considerable manual efforts to accurately specify the input constraints for the SMT solver. 
Regarding the data-driven approaches, even with pre-trained large language models (LLMs), tensor program semantics cannot be fully preserved and thus the translation accuracy only achieves 29.6\%~\cite{li2022competition}, which requires considerable manual efforts for code correction.
In summary, existing transcompilation techniques for tensor programs face challenges such as tremendous manual efforts, limited scalability, or functional incorrectness, as shown in Figure~\ref{fig:intro_comparison}.



To automatically translate tensor programs across deep learning systems with a correctness guarantee, we propose to exploit both the \emph{flexibility of data-driven approaches}, specifically, via LLMs, and the \emph{soundness of symbolic synthesis approaches}. However, the challenge of collaboration between LLM and symbolic synthesis is three-fold: (1) it is difficult to generate highly accurate code via LLM due to the scarcity of tensor programs, (2) it is non-trivial to determine which code snippets should be formally generated by search-based symbolic synthesis, in order to balance the computational costs and achieved accuracy, and (3) performance optimizations are hard to specify either with LLM prompts or synthesis specifications, making the performance of generated code significantly lags behind that of human experts.

%\vspace{-15pt}
\subsection{Our Proposal}
To address the above challenges, we propose to build a novel transcompiler, i.e., \name{}, for translating tensor programs across DLS via both LLM (i.e., GPT-4~\cite{bubeck2023sparks}) and search-based symbolic synthesis. The key of \name{} is that \emph{the program translation is automatically conducted as a series of LLM-assisted transformation passes, where incorrect code snippets are repaired by small-scale symbolic synthesis, and the optimal transformation passes are identified via hierarchical auto-tuning}. The advantages of proposed neural-symbolic synthesis are: (1) the entire transcompilation is decomposed into multiple transformation passes via a chain of LLM prompting, instead of a single prompt in traditional LLM-assisted code generation, to significantly improve the accuracy, (2) the problem size for symbolic synthesis is constrained to a limited scale so that the SMT solver can handle it efficiently, and (3) the performance optimization is achieved by exploring both the parameters and sequences of transformation passes.


\begin{table*}[t]
\vspace{-10pt}
    \caption{Comparison of deep learning systems (Intel DL Boost, NVIDIA GPU, AMD MI, and Cambricon MLU) and their programming interfaces. The characteristics fall into three categories: Parallelism, Memory Hierarchy, and Specialized Intrinsics.}
    \label{tab:my_label}
    \centering
    %\footnotesize
    \scriptsize
    \vspace{-5pt}
\begin{tabular}{ccll}
\toprule
\textbf{Platforms}             & \textbf{Interfaces}     & \multicolumn{1}{c}{\textbf{Categories}}         & \multicolumn{1}{c}{\textbf{Examples}}                                                                              \\ 
\midrule
Intel DL Boost               & \begin{tabular}[c]{@{}c@{}}C with\\ VNNI extensions\end{tabular}        & Specialized Intrinsic                & \texttt{\_mm\_dpbusds\_epi32(...), \_mm512\_dpbusd\_epi32(...)} \\                                                                        
\midrule
\multirow{4}{*}{\begin{tabular}[c]{@{}c@{}}NVIDIA GPU\\ with Tensor Core\end{tabular}}    & \multirow{4}{*}{CUDA C} & Parallelism                                   & \texttt{blockIdx, threadIdx}                                                                                       \\ \cmidrule{3-4} 
                               &                         & \multirow{2}{*}{Memory Hierarchy}             & \texttt{registers, \_\_shared\_\_, \_\_global\_\_}                                                                 \\ \cmidrule{4-4} 
                               &                         &                                               & \begin{tabular}[c]{@{}l@{}}\texttt{matrix\_a}, \texttt{matrix\_b}, and \texttt{accumulator} in Tensor Core fragments\end{tabular} \\ \cmidrule{3-4} 
                               &                         & Specialized Intrinsics                  & \texttt{wmma::mma\_sync(d, a, b, c)}                                                                               \\ \midrule
\multirow{4}{*}{\begin{tabular}[c]{@{}c@{}}AMD MI\\ with Matrix Core\end{tabular}}        & \multirow{4}{*}{HIP}    & Parallelism                                   & \texttt{blockIdx, threadIdx}                                                                                 \\ \cmidrule{3-4} 
                               &                         & \multirow{2}{*}{Memory Hierarchy}             & \texttt{registers, \_\_shared\_\_, \_\_global\_\_}                                                                 \\ \cmidrule{4-4} 
                               &                         &                                               & \begin{tabular}[c]{@{}l@{}}\texttt{matrix\_a}, \texttt{matrix\_b}, and \texttt{accumulator} in Matrix Core fragments\end{tabular} \\ \cmidrule{3-4} 
                               &                         & Specialized Intrinsics                  & \texttt{d = \_\_builtin\_amdgcn\_mfma\_f32\_16x16x4f32(a, b, c, ...)}                                              \\ \midrule
\multirow{6}{*}{Cambricon MLU} & \multirow{6}{*}{BANG C} & \multirow{2}{*}{Parallelism}                  & \texttt{taskId} for task-level parallelism                                                                                                    \\ \cmidrule{4-4} 
                               &                         &                                               & \texttt{clusterId, coreId}    for multi-core parallelism                                                                                     \\ \cmidrule{3-4} 
                               &                         & \multirow{2}{*}{Memory Hierarchy}             & \texttt{registers, \_\_mlu\_shared\_\_, \_\_mlu\_device\_\_}                                                       \\ \cmidrule{4-4} 
                               &                         &                                               & \texttt{\_\_nram\_\_, \_\_wram\_\_}                                                                                \\ \cmidrule{3-4} 
                               &                         & Specialized Intrinsics & \texttt{\_\_bang\_mlp(...)}, \texttt{\_\_bang\_conv(...)}                                                                                                                   \\ \bottomrule
\end{tabular}

\vspace{-10pt}   
\end{table*}


Concretely, by considering the high parallelism, memory hierarchy, and specialized ISA of DLS, the proposed approach consists of three categories of transformation, i.e., \emph{sequentialization/parallelization}, \emph{memory conversion}, and \emph{(de)tensorization}. Each transformation pass is first performed with LLM-based code transformation, then validated with unit-test, and finally repaired with an SMT-solver if necessary. 
Then, a hierarchical auto-tuning approach consisting of the \emph{intra-pass} and \emph{inter-pass} auto-tuning is conducted to improve the performance. The intra-pass auto-tuning uses brute-force search to find the optimal parameters (e.g., tiling sizes) for program transformation, and the inter-pass auto-tuning uses MCTS (Monte Carlo Tree Search)~\cite{6145622} to determine the optimal sequence of passes with maximized performance.


To demonstrate the generality of \name{}, we conduct experiments on 4 different DLS with their interfaces, i.e.,  Intel DL Boost with
VNNI intrinsics, NVIDIA GPU with CUDA C, AMD MI with HIP code, and
Cambricon MLU with BANG C. Experimental results show that \name{} correctly translates different programs at the accuracy of 95\% on average, 
and the performance of translated programs achieves up to $2.0\times$ over \todo{vendor-provided manually-optimized libraries including cuDNN/cuBLAS~\cite{NVCUDNN,NVCUBLAS} and oneDNN~\cite{oneDNN}}.
Moreover, the programming productivity of $2$ representative DLS, i.e., NVIDIA GPU and MLU, is improved by up to $34.3\times$ and $96.0\times$, respectively, via transcompiling legacy tensor programs.

\subsection{Key Contributions}
To our best knowledge, this work is \emph{the first to automatically translate tensor programs of deep learning systems with different programming models}. This paper makes the following contributions:
\begin{itemize}

\item {\textbf{Neural-symbolic program synthesis}}. We propose to use LLMs for generating high-level program sketches via pre-defined meta-prompts and repair the incorrect low-level details through SMT-based symbolic synthesis with limited scale, which achieves \emph{fully automatic program translation with a correctness guarantee}.

\item{\textbf{Hierarchical performance auto-tuning}}. We propose a hierarchical auto-tuning approach, where the \emph{intra-pass} and \emph{inter-pass} auto-tuning are used for exploring pass parameters and pass sequences, respectively, to maximize the performance of synthesized programs.



\item{\textbf{Extensive evaluation}}. We conduct extensive evaluations on 4 different deep learning systems, and experimental results on accuracy, execution performance, and productivity improvement well demonstrate the effectiveness and efficiency of this work.

\end{itemize}